\section{Main Theorems}\label{main-theorems}

We prove the core results connecting leverage to error probability and architectural optimality. The theoretical structure is:

\begin{enumerate}
\item \textbf{DOF-Reliability Isomorphism} (Theorem \ref{thm:dof-reliability}): Maps architecture to reliability theory
\item \textbf{Leverage-Error Tradeoff} (Theorem \ref{thm:leverage-error}): Connects leverage to error probability
\item \textbf{Physical Edit-Energy Floor} (Theorem \ref{thm:physical-energy-floor}): DOF controls minimum edit-energy
\item \textbf{Optimality Criterion} (Theorem \ref{thm:optimal}): Correctness in a fixed capability class
\item \textbf{Composition Stability} (Theorem \ref{thm:composition}): Composition preserves leverage lower bounds
\end{enumerate}

\subsection{Recap: DOF-Reliability Isomorphism}

The core theoretical contribution (stated in Section 1.3) is that DOF maps formally to series system components. This enables all subsequent results by connecting software architecture to the mature mathematical framework of reliability theory.

\textbf{Key properties of the isomorphism:}
\begin{itemize}
\item \textbf{Preserves ordering:} If $\text{DOF}(A_1) < \text{DOF}(A_2)$, then $P_{\text{error}}(A_1) < P_{\text{error}}(A_2)$
\item \textbf{Invertible:} An architecture can be reconstructed from its series system representation
\item \textbf{Compositional:} $\text{DOF}(A_1 \oplus A_2) = \text{DOF}(A_1) + \text{DOF}(A_2)$ (series systems combine)
\end{itemize}

\subsection{The Leverage Maximization Principle}

Given the DOF-Reliability Isomorphism, the following is a \emph{corollary}:

\begin{theorem}[Leverage Maximization Principle]\label{thm:leverage-max}
For any architectural decision with alternatives $A_1, \ldots, A_n$ meeting capability requirements, the optimal choice maximizes leverage:
\[
A^* = \arg\max_{A_i} L(A_i) = \arg\max_{A_i} \frac{|\text{Capabilities}(A_i)|}{\text{DOF}(A_i)}
\]
\end{theorem}
\leanmeta{\LH{L_ad7ca324},
\LH{L_0c281f80}}

\textbf{Note:} This is \emph{not} the central theorem---it is a consequence of the DOF-Reliability Isomorphism. The isomorphism is the deep result; ``maximize leverage'' is the actionable heuristic derived from it.

\subsection{Leverage-Error Tradeoff}

\begin{theorem}[Leverage-Error Tradeoff]\label{thm:leverage-error}
For architectures $A_1, A_2$ with equal capabilities:
\[
\text{Cap}(A_1) = \text{Cap}(A_2) \wedge L(A_1) > L(A_2) \implies P_{\text{error}}(A_1) < P_{\text{error}}(A_2)
\]
\end{theorem}
\leanmeta{\LH{L_5af424e9}}

\begin{proof}
Given: $\text{Cap}(A_1) = \text{Cap}(A_2)$ and $L(A_1) > L(A_2)$.

Since $L(A) = |\text{Cap}(A)|/\text{DOF}(A)$ and capabilities are equal:
\[
\frac{|\text{Cap}(A_1)|}{\text{DOF}(A_1)} > \frac{|\text{Cap}(A_2)|}{\text{DOF}(A_2)}
\]

With $|\text{Cap}(A_1)| = |\text{Cap}(A_2)|$:
\[
\frac{1}{\text{DOF}(A_1)} > \frac{1}{\text{DOF}(A_2)} \implies \text{DOF}(A_1) < \text{DOF}(A_2)
\]

By Corollary \ref{cor:dof-monotone}:
\[
\text{DOF}(A_1) < \text{DOF}(A_2) \implies P_{\text{error}}(A_1) < P_{\text{error}}(A_2)
\]
\end{proof}

\textbf{Corollary:} Maximizing leverage minimizes error probability (for fixed capabilities).

\subsection{Physical Edit-Energy Floor}

\begin{theorem}[Physical Edit-Energy Floor]\label{thm:physical-energy-floor}
Fix a physical edit model $\mathcal{M}$ with per-edit conversion constant $j_{\mathcal{M}}>0$. For any architecture $A$:
\[
E_{\min}(A;\mathcal{M}) = j_{\mathcal{M}}\cdot \mathrm{DOF}(A).
\]
Hence $\mathrm{DOF}(A_1) < \mathrm{DOF}(A_2)$ implies
$E_{\min}(A_1;\mathcal{M}) < E_{\min}(A_2;\mathcal{M})$.
\end{theorem}
\leanmeta{\LH{L_515d1d9d},
\LH{L_5b347895},
\LH{L_05a51b10}}

\begin{proof}
In the mechanized model, modification complexity is definitionally $\mathrm{DOF}$. Multiplying by a positive per-edit conversion constant gives the lower bound and strict monotonicity in DOF.
\end{proof}

\begin{corollary}[Leverage-Energy Monotonicity in a Capability Class]\label{cor:leverage-energy}
For architectures $A_1,A_2$ with equal capabilities, if $L(A_1)>L(A_2)$ then
\[
E_{\min}(A_1;\mathcal{M}) < E_{\min}(A_2;\mathcal{M}),
\]
and the induced energy gap is strictly positive.
\end{corollary}
\leanmeta{\LH{L_9c6279a3},
\LH{L_ebcd21af}}

\begin{theorem}[Budget Feasibility Boundary]\label{thm:physical-budget-boundary}
For physical model $\mathcal{M}$ and budget $B$, implementation feasibility is equivalent to clearing the floor:
\[
E_{\min}(A;\mathcal{M}) \le B \iff \text{feasible}(A,B,\mathcal{M}).
\]
Hence $B < E_{\min}(A;\mathcal{M})$ implies infeasibility. Also, for fixed $B$ and equal capabilities, if $L(A_1)>L(A_2)$ and $A_2$ is feasible under $B$, then $A_1$ is feasible under $B$.
\end{theorem}
\leanmeta{\LH{L_1af867ce},
\LH{L_c4bbc516},
\LH{L_3432eddd}}

\begin{corollary}[Physical Assumption Necessity Witnesses]\label{cor:physical-assumption-necessity}
The physical layer uses two explicit premises for no-go style infeasibility claims: positive per-edit conversion and an external budget bound. Dropping positivity admits a zero-floor countermodel. Dropping the external budget-bound premise blocks infeasibility conclusions because a feasible budget witness always exists.
\end{corollary}
\leanmeta{\LH{L_a2e95cfd},
\LH{L_f3216891},
\LH{L_e0961ce8},
\LH{L_1fb0826b}}

\subsection{Metaprogramming Dominance}

\begin{theorem}[Metaprogramming Dominance]\label{thm:metaprog}
For any $N \in \mathbb{N}$, there exists a metaprogramming architecture $M$ with $\text{DOF}(M) = 1$ and $L(M) \geq N$.
\end{theorem}
\leanmeta{\LH{L_c022ff39},
\LH{L_5250a1d9}}

\begin{proof}
Let $M$ be a metaprogramming architecture with a single source definition ($\text{DOF}(M) = 1$) and at least $N$ derived capabilities. Then $L(M) = |\text{Cap}(M)|/1 \geq N$.

\textbf{Modeling note.} In any fixed finite model, $L(M)$ is a positive rational. The colloquial claim ``leverage $= \infty$'' means: for any bound $N$, the architecture can be extended to achieve $L \geq N$ while keeping $\text{DOF} = 1$. The Lean formalization proves the for-all-$N$ version (\LH{L_c022ff39}).
\end{proof}

\subsection{Architectural Decision Criterion}

\begin{theorem}[Optimal Architecture]\label{thm:optimal}
Given requirements $R$, architecture $A^*$ is optimal in its capability class if and only if:
\begin{enumerate}
\item $\text{Cap}(A^*) \supseteq R$ (feasibility)
\item $\forall A'$ with $\text{Cap}(A') = \text{Cap}(A^*)$ and $\text{Cap}(A') \supseteq R$: $L(A^*) \geq L(A')$ (maximality in fixed capability class)
\end{enumerate}
\end{theorem}
\leanmeta{\LH{L_a42f986c},
\LH{L_28e01ed8}}

\begin{proof}
($\Leftarrow$) Suppose $A^*$ satisfies (1) and (2). Then $A^*$ is feasible and has maximum leverage among feasible architectures in the same capability class. By Theorem \ref{thm:leverage-error}, this minimizes error probability in that class.

($\Rightarrow$) If (1) fails, $A^*$ is infeasible. If (2) fails, there exists $A'$ in the same capability class with $L(A') > L(A^*)$, so $P_{\text{error}}(A') < P_{\text{error}}(A^*)$ by Theorem \ref{thm:leverage-error}, contradicting optimality.
\end{proof}

\textbf{Decision Procedure:}
\begin{enumerate}
\item Enumerate candidate architectures $\{A_1, \ldots, A_n\}$
\item Filter: Keep only $A_i$ with $\text{Cap}(A_i) \supseteq R$
\item Optimize: Choose $A^* = \arg\max_i L(A_i)$
\end{enumerate}

\subsection{Leverage Composition}

\begin{theorem}[Leverage Composition]\label{thm:composition}
For modular architecture $A = A_1 \oplus A_2$ with disjoint components:
\begin{enumerate}
\item $\text{DOF}(A) = \text{DOF}(A_1) + \text{DOF}(A_2)$
\item if $L(A_1)\ge 1$ and $L(A_2)\ge 1$, then $L(A)\ge 1$
\end{enumerate}
\end{theorem}
\leanmeta{\LH{L_1bc009bb},
\LH{L_509b7c92},
\LH{L_b8e54251}}

\begin{proof}
(1) By Proposition \ref{prop:dof-additive}.

(2) If $L(A_i)\ge 1$, then $|\text{Cap}(A_i)|\ge \text{DOF}(A_i)$ for $i=1,2$. Summing both inequalities gives
\[
|\text{Cap}(A_1)|+|\text{Cap}(A_2)|\ge \text{DOF}(A_1)+\text{DOF}(A_2),
\]
which is exactly $L(A)\ge 1$ under additive composition.
\end{proof}

\textbf{Interpretation:} Composition preserves a baseline leverage floor under the stated assumptions.

\begin{remark}[Composition Breaks SSOT]\label{rem:composition-breaks-ssot}
Theorem~\ref{thm:composition} preserves leverage \emph{floors}, but composing two SSOT architectures breaks SSOT: if $\mathrm{DOF}(A_1) = \mathrm{DOF}(A_2) = 1$, then $\mathrm{DOF}(A_1 \oplus A_2) = 2$. This is formalized as \texttt{compose\_breaks\_ssot} in \texttt{Leverage/BridgeToDQ.lean}. The composition tax is unavoidable: distributed systems consisting of $k$ independent SSOT components have $\mathrm{DOF} = k$, placing them in the coNP-hard regime (srank $= k > 1$) with mandatory thermodynamic cost per decision cycle.
\end{remark}

\subsection{Formalization}

Main optimization theorems are formalized in \texttt{Leverage/Theorems.lean}; physical edit-energy theorems are formalized in \texttt{Leverage/Physical.lean}:
\begin{itemize}
\item \LH{L_22638fcd}: Theorem \ref{thm:leverage-error}
\item \LH{L_9dc717da} and
  \LH{L_bc94b2f8}:
  Theorem \ref{thm:physical-energy-floor}, Corollary \ref{cor:leverage-energy}
\item \LH{L_53bac39c},
  \LH{L_dd0fab2d}, and
  \LH{L_e5602ce9}:
  Theorem \ref{thm:physical-budget-boundary},
  Corollary \ref{cor:physical-assumption-necessity}
\item \LH{L_b85b0db1}: Theorem \ref{thm:metaprog}
\item \LH{L_bb990fcf} and \LH{L_7254810e}:
  Theorem \ref{thm:optimal}
\item \LH{L_a4e66dad},
  \LH{L_f7a7de2d}, and
  \LH{L_0869dd64}:
  Theorem \ref{thm:composition}
\item \texttt{Leverage/Physical.lean}: physical edit-energy floor theorems
\end{itemize}

\subsection{Integration with Dependencies}

The leverage framework provides the unifying theory for results proved in \texttt{AbstractClassSystem} and \texttt{Ssot}:

\begin{theorem}[\texttt{AbstractClassSystem} as Leverage Instance]\label{thm:paper1-integration}
Nominal typing dominance from \texttt{AbstractClassSystem}~\cite{paper1_typing_discipline} is an instance of leverage maximization:
\begin{itemize}
\item Nominal typing adds 4 B-dependent capabilities over duck typing
\item DOF remains fixed in the mechanized typing instance
\item Therefore nominal typing has strictly higher leverage
\end{itemize}
\end{theorem}
\leanmeta{\LH{L_9f6625e0},
\LH{L_4372537d}}

\begin{theorem}[\texttt{Ssot} as Leverage Instance]\label{thm:paper2-integration}
SSOT dominance from \texttt{Ssot}~\cite{paper2_ssot} is an instance of leverage maximization:
\begin{itemize}
\item SSOT fixes DOF at 1 for a structural fact
\item Non-SSOT replication yields DOF $= n$ for the same fact
\item Therefore leverage improves by factor $n$
\end{itemize}
\end{theorem}
\leanmeta{\LH{L_3c6b4088},
\LH{L_48108f64}}

These integration theorems are formalized in:
\begin{itemize}
\item \nolinkurl{Leverage/Typing.lean}
\item \nolinkurl{Leverage/SSOT.lean}
\end{itemize}
